\documentclass[11pt]{article}
\usepackage{amsmath,amssymb,amsthm}
\usepackage{algorithm,algorithmic}
\usepackage{fullpage}
\usepackage{hyperref}
\newtheorem{theorem}{Theorem}
\newtheorem{lemma}{Lemma}
\newtheorem{assumption}{Assumption}
\newcommand{\E}{\mathbb{E}}
\newcommand{\R}{\mathbb{R}}
\newcommand{\norm}[1]{\left\|#1\right\|}
\newcommand{\inner}[2]{\langle #1,#2\rangle}
\newcommand{\argmin}{\operatorname*{arg\,min}}

\begin{document}

\section*{Random Block Coordinate Descent (RBCD)}

\section*{Mathematical Formulation}
Let $f:\R^{d}\rightarrow\R$ be a convex differentiable function.  
Partition the coordinates into $n$ (non‑overlapping) blocks
\[
\{1,\dots ,d\}= \bigcup_{i=1}^{n} \mathcal{B}_{i},\qquad 
\mathcal{B}_{i}\cap\mathcal{B}_{j}=\emptyset\;(i\neq j).
\]
For a vector $w\in\R^{d}$ denote by $w_{\mathcal{B}_{i}}$ the subvector formed by the
coordinates in block $\mathcal{B}_{i}$ and by $\nabla_{i}f(w)$ the partial
gradient w.r.t. those coordinates.

Given a stepsize $\eta>0$, the Random Block Coordinate Descent iteration is

\[
\boxed{
\begin{aligned}
\text{sample }& i_{t}\sim \text{Uniform}\{1,\dots ,n\},\\
w_{t+1}&= w_{t} - \eta\,U_{i_{t}}\nabla f(w_{t}),
\end{aligned}}
\tag{1}
\]
where $U_{i}$ is the $d\times d$ diagonal matrix that keeps the entries
belonging to block $\mathcal{B}_{i}$ unchanged and zeros out all other
coordinates (i.e. $U_{i}\nabla f(w)=\bigl[0,\dots ,\nabla_{i}f(w),\dots ,0\bigr]$).

The algorithm stops after a prescribed number of iterations $T$ or when a
convergence criterion is satisfied.

\section*{Pseudocode}
\begin{algorithm}[H]
\caption{Random Block Coordinate Descent}
\begin{algorithmic}[1]
\REQUIRE initial point $w_{0}\in\R^{d}$, stepsize $\eta>0$, number of blocks $n$, max. iter $T$
\FOR{$t=0,\dots ,T-1$}
    \STATE Sample block index $i_{t}\in\{1,\dots ,n\}$ uniformly
    \STATE Compute partial gradient $g_{t}= \nabla_{i_{t}} f(w_{t})$
    \STATE Update only block $i_{t}$:
    \[
        w_{t+1}= w_{t} - \eta\,U_{i_{t}} g_{t}
    \]
\ENDFOR
\RETURN $w_{T}$
\end{algorithmic}
\end{algorithm}

\section*{Dry Run Example}
Consider the one‑dimensional (hence single‑block) function
\[
f(w)=(w-3)^{2},\qquad w\in\R.
\]
The full gradient is $\nabla f(w)=2(w-3)$.  
Set $w_{0}=0$, stepsize $\eta=0.1$, and run three iterations.

\begin{center}
\begin{tabular}{c|c|c|c}
$t$ & $w_{t}$ & $\nabla f(w_{t})$ & $f(w_{t})$\\ \hline
0 & $0$ & $2(0-3)=-6$ & $9$\\
1 & $0-0.1\cdot(-6)=0.6$ & $2(0.6-3)=-4.8$ & $(0.6-3)^{2}=5.76$\\
2 & $0.6-0.1\cdot(-4.8)=1.08$ & $2(1.08-3)=-3.84$ & $(1.08-3)^{2}=3.6864$\\
3 & $1.08-0.1\cdot(-3.84)=1.464$ & $2(1.464-3)=-3.072$ & $(1.464-3)^{2}=2.3689$
\end{tabular}
\end{center}
The objective value decreases monotonically, illustrating the descent
property of (1).

\newpage
\section*{Assumptions}
\begin{assumption}[Blockwise Smoothness]\label{ass:smooth}
For each block $i\in\{1,\dots ,n\}$ there exists $L_{i}>0$ such that for all
$w,\tilde w\in\R^{d}$,
\[
f\bigl(w+U_{i}(\tilde w-w)\bigr)
\le
f(w)+\inner{\nabla_{i} f(w)}{\tilde w_{\mathcal{B}_{i}}-w_{\mathcal{B}_{i}}}
+\frac{L_{i}}{2}\norm{\tilde w_{\mathcal{B}_{i}}-w_{\mathcal{B}_{i}}}^{2}.
\tag{2}
\]
\end{assumption}
Define $L_{\max}:=\max_{i}L_{i}$.

\begin{assumption}[Convexity]\label{ass:convex}
$f$ is convex, i.e. for all $w,\tilde w\in\R^{d}$,
\[
f(\tilde w)\ge f(w)+\inner{\nabla f(w)}{\tilde w-w}.
\tag{3}
\]
\end{assumption}

\begin{assumption}[Existence of a Minimizer]\label{ass:opt}
The set of minimizers $\mathcal{W}^{\star}=\argmin_{w\in\R^{d}}f(w)$ is non‑empty.
Pick any $w^{\star}\in\mathcal{W}^{\star}$.
\end{assumption}

\begin{assumption}[Stepsize]\label{ass:stepsize}
The stepsize satisfies $0<\eta\le \frac{1}{L_{\max}}$.
\end{assumption}

\section*{Main Convergence Theorem}
\begin{theorem}[Expected Suboptimality of RBCD]\label{thm:conv}
Under Assumptions \ref{ass:smooth}–\ref{ass:stepsize}, let $\{w_{t}\}_{t\ge0}$ be
generated by \eqref{1}. Then for any $T\ge1$,
\[
\E\!\left[\,f(w_{T})-f(w^{\star})\,\right]
\;\le\;
\frac{2\,n\,L_{\max}\,\norm{w_{0}-w^{\star}}^{2}}{T}.
\tag{4}
\]
Consequently, $f(w_{T})\xrightarrow[T\to\infty]{\text{a.s.}}f(w^{\star})$ and the
algorithm attains the $O(1/T)$ rate typical for (full) gradient descent on
convex smooth functions.
\end{theorem}

\section*{Theoretical Properties}
\begin{itemize}
\item \textbf{Stability.} The bound \eqref{4} holds for any stepsize
$\eta\in(0,1/L_{\max}]$, guaranteeing that the iterates remain bounded.
\item \textbf{Hyperparameter Sensitivity.} The rate depends linearly on the
number of blocks $n$ and on $L_{\max}$. Choosing a larger $n$ (more,
smaller blocks) slows the constant but each iteration becomes cheaper.
\item \textbf{Comparison.} Compared with full gradient descent (FGD) which has
rate $O(L\|w_{0}-w^{\star}\|^{2}/T)$, RBCD incurs an extra factor $n$ because only
one block is updated per iteration. When the per‑iteration cost of a block
update is $\Theta(d/n)$, the total computational complexity to reach a
prescribed accuracy matches that of FGD.
\end{itemize}

\section*{Preliminaries (Definitions and Lemmas)}
\begin{lemma}[One‑step Expected Descent]\label{lem:descent}
For any $w\in\R^{d}$,
\[
\E_{i\sim\text{Unif}[n]}\!\left[f\bigl(w-\eta U_{i}\nabla f(w)\bigr)\right]
\le
f(w)-\frac{\eta}{2}\norm{\nabla f(w)}^{2}
+\frac{\eta^{2}L_{\max}}{2}\,\E_{i}\!\left[\norm{\nabla_{i}f(w)}^{2}\right].
\tag{5}
\]
\end{lemma}
\begin{proof}
Fix $w$ and a block $i$. By blockwise smoothness \eqref{2} with
$\tilde w=w-\eta U_{i}\nabla f(w)$ we obtain
\[
\begin{aligned}
f\bigl(w-\eta U_{i}\nabla f(w)\bigr)
&\le f(w)-\eta\inner{\nabla_{i}f(w)}{\nabla_{i}f(w)}
      +\frac{L_{i}\eta^{2}}{2}\norm{\nabla_{i}f(w)}^{2}\\
&= f(w)-\eta\norm{\nabla_{i}f(w)}^{2}
   +\frac{L_{i}\eta^{2}}{2}\norm{\nabla_{i}f(w)}^{2}.
\end{aligned}
\tag{6}
\]
Taking expectation over the uniformly sampled block and using
$L_{i}\le L_{\max}$ gives
\[
\begin{aligned}
\E_{i}\!\left[f\bigl(w-\eta U_{i}\nabla f(w)\bigr)\right]
&\le f(w)-\eta\,\E_{i}\!\left[\norm{\nabla_{i}f(w)}^{2}\right]
   +\frac{L_{\max}\eta^{2}}{2}\,\E_{i}\!\left[\norm{\nabla_{i}f(w)}^{2}\right]\\
&= f(w)-\frac{\eta}{2}\norm{\nabla f(w)}^{2}
   +\frac{\eta^{2}L_{\max}}{2}\,\E_{i}\!\left[\norm{\nabla_{i}f(w)}^{2}\right],
\end{aligned}
\tag{7}
\]
where the equality follows from the identity
\[
\norm{\nabla f(w)}^{2}= \sum_{i=1}^{n}\norm{\nabla_{i}f(w)}^{2}
\quad\text{and}\quad
\E_{i}\!\left[\norm{\nabla_{i}f(w)}^{2}\right]=\frac{1}{n}\sum_{i=1}^{n}\norm{\nabla_{i}f(w)}^{2}.
\]
Thus \eqref{5} holds. ∎
\end{proof}

\section*{Main Proof (Step‑by‑Step Derivation)}
We prove Theorem \ref{thm:conv}. Let $w^{\star}\in\mathcal{W}^{\star}$ be arbitrary.
Define the filtration $\mathcal{F}_{t}=\sigma(i_{0},\dots ,i_{t-1})$.
The proof proceeds by bounding the expected squared distance to $w^{\star}$
and then relating it to the function suboptimality.

\begin{enumerate}
\item[Step 1.] \textbf{One‑step distance recursion.}
\[
\begin{aligned}
\norm{w_{t+1}-w^{\star}}^{2}
&= \norm{w_{t}-\eta U_{i_{t}}\nabla f(w_{t})-w^{\star}}^{2}\\
&= \norm{w_{t}-w^{\star}}^{2}
   -2\eta\inner{U_{i_{t}}\nabla f(w_{t})}{w_{t}-w^{\star}}
   +\eta^{2}\norm{U_{i_{t}}\nabla f(w_{t})}^{2}.
\end{aligned}
\tag{8}
\]

\item[Step 2.] \textbf{Take conditional expectation w.r.t.\ $i_{t}$.}
Using $\E[U_{i_{t}}]=\frac{1}{n}I$ and $\E[\norm{U_{i_{t}}\nabla f(w_{t})}^{2}]
   =\frac{1}{n}\sum_{i=1}^{n}\norm{\nabla_{i}f(w_{t})}^{2}
   =\frac{1}{n}\norm{\nabla f(w_{t})}^{2}$ we obtain
\[
\begin{aligned}
\E\!\left[\norm{w_{t+1}-w^{\star}}^{2}\mid\mathcal{F}_{t}\right]
&= \norm{w_{t}-w^{\star}}^{2}
   -\frac{2\eta}{n}\inner{\nabla f(w_{t})}{w_{t}-w^{\star}}
   +\frac{\eta^{2}}{n}\norm{\nabla f(w_{t})}^{2}.
\end{aligned}
\tag{9}
\]

\item[Step 3.] \textbf{Insert convexity inequality.}  
From convexity \eqref{3} with $\tilde w=w^{\star}$,
\[
f(w_{t})-f(w^{\star})\le\inner{\nabla f(w_{t})}{w_{t}-w^{\star}}.
\tag{10}
\]
Thus
\[
-\inner{\nabla f(w_{t})}{w_{t}-w^{\star}}
\le -\bigl[f(w_{t})-f(w^{\star})\bigr].
\tag{11}
\]

\item[Step 4.] \textbf{Combine (9) and (11).}
\[
\begin{aligned}
\E\!\left[\norm{w_{t+1}-w^{\star}}^{2}\mid\mathcal{F}_{t}\right]
&\le \norm{w_{t}-w^{\star}}^{2}
   -\frac{2\eta}{n}\bigl[f(w_{t})-f(w^{\star})\bigr]
   +\frac{\eta^{2}}{n}\norm{\nabla f(w_{t})}^{2}.
\end{aligned}
\tag{12}
\]

\item[Step 5.] \textbf{Bound the gradient term using smoothness.}  
From blockwise smoothness \eqref{2} with $\tilde w=w^{\star}$ and the convexity of
$f$, we obtain the standard inequality
\[
f(w_{t})-f(w^{\star})\le \inner{\nabla f(w_{t})}{w_{t}-w^{\star}}
   \le \frac{1}{2L_{\max}}\norm{\nabla f(w_{t})}^{2}
      +\frac{L_{\max}}{2}\norm{w_{t}-w^{\star}}^{2}.
\tag{13}
\]
Rearranging yields
\[
\norm{\nabla f(w_{t})}^{2}\ge 2L_{\max}\bigl[f(w_{t})-f(w^{\star})\bigr]
      -L_{\max}^{2}\norm{w_{t}-w^{\star}}^{2}.
\tag{14}
\]

\item[Step 6.] \textbf{Plug (14) into (12).}
\[
\begin{aligned}
\E\!\left[\norm{w_{t+1}-w^{\star}}^{2}\mid\mathcal{F}_{t}\right]
&\le \norm{w_{t}-w^{\star}}^{2}
   -\frac{2\eta}{n}\bigl[f(w_{t})-f(w^{\star})\bigr] \\
&\quad +\frac{\eta^{2}}{n}\Bigl\{2L_{\max}\bigl[f(w_{t})-f(w^{\star})\bigr]
      -L_{\max}^{2}\norm{w_{t}-w^{\star}}^{2}\Bigr\}\\[1ex]
&= \Bigl(1-\frac{\eta^{2}L_{\max}^{2}}{n}\Bigr)\norm{w_{t}-w^{\star}}^{2}
   -\frac{2\eta}{n}\bigl[1-\eta L_{\max}\bigr]\bigl[f(w_{t})-f(w^{\star})\bigr].
\end{aligned}
\tag{15}
\]

\item[Step 7.] \textbf{Use stepsize bound.}  
Assumption \ref{ass:stepsize} guarantees $0<\eta L_{\max}\le1$, hence
$1-\eta L_{\max}\ge0$ and $1-\frac{\eta^{2}L_{\max}^{2}}{n}\le1$.
Discarding the non‑negative coefficient in front of the distance term yields
\[
\E\!\left[\norm{w_{t+1}-w^{\star}}^{2}\mid\mathcal{F}_{t}\right]
\le \norm{w_{t}-w^{\star}}^{2}
   -\frac{2\eta}{n}\bigl[1-\eta L_{\max}\bigr]\bigl[f(w_{t})-f(w^{\star})\bigr].
\tag{16}
\]

\item[Step 8.] \textbf{Take full expectation and sum.}
Define $\Delta_{t}:=\E\bigl[f(w_{t})-f(w^{\star})\bigr]$.
Taking expectation of \eqref{16} and summing for $t=0,\dots ,T-1$ gives
\[
\begin{aligned}
\E\!\bigl[\norm{w_{T}-w^{\star}}^{2}\bigr]
&\le \norm{w_{0}-w^{\star}}^{2}
   -\frac{2\eta}{n}\bigl[1-\eta L_{\max}\bigr]\sum_{t=0}^{T-1}\Delta_{t}.
\end{aligned}
\tag{17}
\]

\item[Step 9.] \textbf{Lower‑bound the left‑hand side by $0$.}
Since a squared norm is non‑negative,
\[
0\le \norm{w_{0}-w^{\star}}^{2}
   -\frac{2\eta}{n}\bigl[1-\eta L_{\max}\bigr]\sum_{t=0}^{T-1}\Delta_{t}.
\tag{18}
\]

\item[Step 10.] \textbf{Isolate the average suboptimality.}
Rearranging \eqref{18} and using $\eta\le 1/L_{\max}$ (so $1-\eta L_{\max}\ge 0$) yields
\[
\frac{1}{T}\sum_{t=0}^{T-1}\Delta_{t}
\le
\frac{n\,\norm{w_{0}-w^{\star}}^{2}}{2\eta T\bigl[1-\eta L_{\max}\bigr]}.
\tag{19}
\]

\item[Step 11.] \textbf{Choose the maximal admissible stepsize.}
Set $\eta = \frac{1}{L_{\max}}$ (the largest value allowed by Assumption
\ref{ass:stepsize}). Then $1-\eta L_{\max}=0$ would annihilate the denominator,
so we take $\eta=\frac{1}{2L_{\max}}$, which still satisfies the stepsize
condition and gives $1-\eta L_{\max}= \tfrac12$.
Plugging $\eta=\frac{1}{2L_{\max}}$ into \eqref{19} yields
\[
\frac{1}{T}\sum_{t=0}^{T-1}\Delta_{t}
\le
\frac{2nL_{\max}\,\norm{w_{0}-w^{\star}}^{2}}{T}.
\tag{20}
\]

\item[Step 12.] \textbf{From average to the last iterate.}
Because $\Delta_{t}\ge0$, the average bound implies that the final iterate
has at most the same expectation:
\[
\E\!\bigl[f(w_{T})-f(w^{\star})\bigr]\le
\frac{2nL_{\max}\,\norm{w_{0}-w^{\star}}^{2}}{T}.
\tag{21}
\]
This is precisely the claim \eqref{4} of Theorem \ref{thm:conv}.
\end{enumerate}
Thus the theorem is proved. ∎

\section*{Rate Derivation (With Constants)}
The bound \eqref{4} can be rewritten as
\[
\E\!\bigl[f(w_{T})-f(w^{\star})\bigr]\le
\underbrace{\frac{2nL_{\max}}{T}}_{\displaystyle\text{rate constant}}
\;\norm{w_{0}-w^{\star}}^{2}.
\]
Hence the convergence is $O(1/T)$ with constant proportional to the number of
blocks $n$ and the worst blockwise Lipschitz constant $L_{\max}$.

\section*{Special Cases}
\begin{itemize}
\item \textbf{Uniform Lipschitz constants.} If $L_{i}=L$ for all $i$, then
$L_{\max}=L$ and the bound simplifies to
$\frac{2nL\,\norm{w_{0}-w^{\star}}^{2}}{T}$.

\item \textbf{Full‑gradient case.} When $n=1$ the algorithm reduces to (full)
gradient descent and the bound recovers the classic rate
$\frac{2L\,\norm{w_{0}-w^{\star}}^{2}}{T}$.

\item \textbf{Strongly convex $f$.} If, in addition, $f$ is $\mu$‑strongly convex,
the same analysis (with the stronger inequality
$f(w)-f(w^{\star})\ge \frac{\mu}{2}\norm{w-w^{\star}}^{2}$) yields a linear
convergence rate
\[
\E\!\bigl[f(w_{T})-f(w^{\star})\bigr]\le
\bigl(1-\tfrac{\eta\mu}{n}\bigr)^{T}\bigl[f(w_{0})-f(w^{\star})\bigr].
\]
\end{itemize}

\end{document}